\sscp{\tsc{\ili{Gorontalo}-\ili{Mongondow} (Greater \ili{Central} Philippine)}}{C-09-03}
The following languages are all Austronesian,
belonging to the \tsc{\ili{Gorontalo}-Mongondow} branch of the \tsc{Greater
\ili{Central} Philippine} subgroup, but spoken in northern Sulawesi.
\citet{Blust2019} further classifies \tsc{Greater \ili{Central} Philippine}
as a branch of his \tsc{Philippine} subgroup.

\noindent{\wg\st{6pt}\tblr{llllll}
\lsptoprule
%Language	&	ISO	&	Term	&	Gloss	&	Etymon	&	Source	\\	\midrule
%\endfirsthead
Language	&	ISO	&	Term	&	Gloss	&	Etymon	&	Source	\\	\midrule
%\endhead
\ili{Bolango}	&	bld	&	ʔuse	&	bear cuscus	&	**kusay	&	\cite[322]{Usup1986}	\\
\ili{Buol}	&	blf	&	kute	&	bear cuscus	&	**kusay	&	\cite[322]{Usup1986}	\\
\ili{Bintauna}	&	bne	&	ʔuse	&	bear cuscus	&	**kusay	&	\cite[322]{Usup1986}	\\
\ili{Gorontalo}	&	gor	&	bubudu	&	cuscus	&	\it{misc.}	&	\cite[437]{Adriani1902}	\\
\ili{Kaidipang}	&	kzp	&	kuse	&	bear cuscus	&	**kusay	&	\cite[322]{Usup1986}	\\
\ili{Lolak}	&	llq	&	kuse	&	bear cuscus	&	**kusay	&	\cite[322]{Usup1986}	\\
\ili{Mongondow}	&	mog	&	kutoy	&	bear cuscus	&	**kusay	&	\cite[322]{Usup1986}	\\
\ili{Ponosakan}	&	pns	&	kusoy	&	bear cuscus	&	**kusay	&	\cite[322]{Usup1986}	\\
\ili{Suwawa}	&	swu	&	ʔute	&	bear cuscus	&	**kusay	&	\cite[322]{Usup1986}	\\
\lspbottomrule
\end{tabular}\mbox{}\\}

%\noindent{\wg\st{2pt}\begin{tabularx}{\linewidth}
%{lllll>{\raggedright\arraybackslash}X}
%\lsptoprule
%%Language	&	ISO	&	Term	&	Gloss	&	Etymon	&	Source	\\	\midrule
%%\endfirsthead
%Language	&	ISO	&	Term	&	Gloss	&	Etymon	&	Source	\\	\midrule
%\endhead
%
%\lspbottomrule
%\end{tabularx}}

%\noindent{\wg\st{2pt}\begin{xltabular}{\linewidth}
%{lllll>{\raggedright\arraybackslash}X}
%\lsptoprule
%%Language	&	ISO	&	Term	&	Gloss	&	Etymon	&	Source	\\	\midrule
%%\endfirsthead
%Language	&	ISO	&	Term	&	Gloss	&	Etymon	&	Source	\\	\midrule
%\endhead
%
%\lspbottomrule
%\end{xltabular}}